\chapter{Post-Quantum Cryptography}

\providecommand{\ket}[1]{\lvert #1 \rangle}

Every public-key scheme in this book so far rests on one of two hard problems: factoring a large integer (RSA) or computing a discrete logarithm in some group (Diffie-Hellman, ElGamal, DSA, and all of elliptic-curve cryptography). Both have served us well for half a century. Both are broken --- completely, not marginally --- the day someone builds a large, fault-tolerant quantum computer. Not weakened the way DES was weakened by faster machines, where you respond by adding key bits. Broken in a way that no key size can fix, because the quantum algorithm that kills them runs in time polynomial in the key length.

That is a strange situation. The machine that breaks RSA-2048 does not exist yet, and credible people disagree about whether it will exist in ten years or thirty or ever. But the data you encrypt today can be copied today and stored, and decrypted the day the machine arrives. An adversary who is patient --- a nation-state archiving traffic, say --- is already collecting. So the question ``when will quantum computers break RSA'' is the wrong one for a defender. The right question is ``how long does my secret need to stay secret, and how long will migration take,'' and for a great deal of data the answer makes the threat present-tense.

This chapter explains why quantum computers menace exactly the schemes they menace (and, just as importantly, why they leave AES and SHA-256 mostly intact), gives you the intuition behind Shor's and Grover's algorithms without assuming you know quantum mechanics, and then builds up the replacements: lattice problems and the LWE assumption, the Kyber and Dilithium schemes that NIST standardized, and hash-based signatures that lean only on the one primitive we trust most. We will trace small numbers by hand, break a toy lattice scheme to see where the security actually lives, and close with the practical questions of migration, hybrid key exchange, and timelines.

\chapterline

\section{What a quantum computer is, in just enough detail}
\label{ch35:qc-model}

You do not need the physics. You need a computational model precise enough to see why factoring falls and AES does not. Here it is.

A classical bit is $0$ or $1$. A \emph{qubit} has a state described by two complex numbers $\alpha$ and $\beta$ with $\abs{\alpha}^2 + \abs{\beta}^2 = 1$, written
\[
  \ket{\psi} = \alpha\ket{0} + \beta\ket{1}.
\]
The notation $\ket{0}$ is just a name for a basis vector; think of $\ket{0}=\binom{1}{0}$ and $\ket{1}=\binom{0}{1}$. The qubit is a unit vector in $\C^2$. When you \emph{measure} it, you get $0$ with probability $\abs{\alpha}^2$ and $1$ with probability $\abs{\beta}^2$, and the state collapses to whatever you saw. The numbers $\alpha,\beta$ are \emph{amplitudes}, not probabilities; the essential difference is that amplitudes can be negative (or complex), so two contributions to the same outcome can \emph{cancel}. That cancellation---interference---provides the foundation for quantum algorithmic speedups.

With $n$ qubits the state is a unit vector in $\C^{2^n}$: one amplitude for each of the $2^n$ classical strings. A computation is a sequence of \emph{gates}, each a unitary matrix (a length-preserving linear map) acting on a few qubits, followed by a measurement. The catch that kills the naive ``it tries all $2^n$ answers at once'' story: you only get to read out one string, sampled according to $\abs{\text{amplitude}}^2$. A quantum algorithm achieves a speedup only if it arranges, through interference, for the amplitudes of \emph{wrong} answers to cancel and the amplitudes of \emph{right} answers to reinforce, so that the single measurement at the end lands on a useful outcome with high probability.

\begin{intuitionbox}
Picture $2^n$ little arrows (the amplitudes), one per possible answer. A classical randomized algorithm can only make these arrows shorter or longer --- probabilities are non-negative. A quantum algorithm can also \emph{rotate} them, so that arrows pointing in opposite directions sum to nearly zero. Designing a quantum algorithm means choosing unitary steps so that the wrong answers point in opposing directions and cancel out, while the right answers align and reinforce one another. Shor's algorithm achieves this interference pattern using the \emph{period} of a function; Grover's uses a slower, more general sequence of reflections that gradually rotates the state vector toward the target.
\end{intuitionbox}

\begin{historybox}
The idea that quantum systems might compute things classical machines cannot goes back to Richard Feynman (1981/82), who pointed out that simulating quantum physics on a classical computer seems to cost exponential time, and David Deutsch (1985), who formalized the quantum Turing machine. For a decade it was a curiosity. Then in 1994 Peter Shor, at Bell Labs, found polynomial-time quantum algorithms for factoring and discrete log, and the field stopped being a curiosity overnight: the world's deployed public-key cryptography had an expiration date stamped on it. Lov Grover's search algorithm followed in 1996. To build quantum-resistant replacements, Mikl{\'o}s Ajtai connected average-case cryptographic hardness to worst-case lattice problems in 1996, and Oded Regev introduced Learning With Errors (LWE) in 2005. From 2016 to 2024 NIST ran an open competition to standardize post-quantum cryptography: CRYSTALS-Kyber and CRYSTALS-Dilithium (``CRYSTALS'' = Cryptographic Suite for Algebraic Lattices) were selected in 2022 and standardized in 2024 as ML-KEM (FIPS 203) and ML-DSA (FIPS 204), the stateless hash-based SPHINCS+ became SLH-DSA (FIPS 205), and the lattice signature Falcon was selected as FN-DSA (draft FIPS 206).
\end{historybox}

\section{Shor's algorithm: why factoring and discrete log fall}
\label{ch35:shor}

RSA's security rests on the gap between multiplying (easy) and factoring (believed hard classically: the best known algorithm, the general number field sieve, runs in time roughly $\exp\!\big(c\,(\log N)^{1/3} \allowbreak (\log\log N)^{2/3}\big)$, sub-exponential but super-polynomial). Shor's algorithm factors in time polynomial in $\log N$. To see how, we reduce factoring to \emph{period-finding}, and then quantum mechanics does the period-finding fast.

\subsection{Factoring reduces to finding the order of an element}
\label{ch35:order-reduction}

Pick a random $a$ with $1 < a < N$ and $\gcd(a,N)=1$ (if the gcd is not $1$ you already found a factor --- lucky). Consider the function
\[
  f(x) = a^x \bmod N.
\]
Because $a$ is a unit mod $N$, this function is periodic: there is a smallest $r>0$, the \emph{order} of $a$, with $a^r \equiv 1 \pmod N$, and then $f(x+r)=f(x)$ for all $x$. Suppose we can find $r$. Two pieces of luck, each with constant probability over the choice of $a$, finish the job:

\begin{enumerate}
  \item $r$ is even. Then $a^{r/2}$ is a square root of $1$ mod $N$, so $N \mid (a^{r/2}-1)(a^{r/2}+1)$.
  \item $a^{r/2} \not\equiv -1 \pmod N$ (and of course $a^{r/2}\not\equiv 1$, since $r$ is the \emph{smallest} period). Then $a^{r/2}\pm 1$ are nontrivial: neither is a multiple of $N$, yet their product is. So $\gcd(a^{r/2}-1,\,N)$ is a nontrivial factor of $N$, computable in an instant by Euclid.
\end{enumerate}

One can show that for $N$ with at least two distinct odd prime factors, a random $a$ satisfies both conditions with probability at least $1/2$. So factoring reduces to finding the order $r$ of a random element. Everything classical here is cheap. The only hard step is finding $r$, and that is where the quantum computer comes in.

\begin{example}[Factoring $15$ by hand]
\label{ch35:ex-15}
Take $N=15$, $a=7$. Compute powers mod $15$: $7^1=7,\ 7^2=49\equiv 4,\ 7^3\equiv 28\equiv 13,\ 7^4\equiv 91\equiv 1$. So the order is $r=4$, which is even. Then $a^{r/2}=7^2=49\equiv 4 \pmod{15}$, and $4\not\equiv \pm1$. Now
\[
  \gcd(4-1,\,15)=\gcd(3,15)=3,\qquad \gcd(4+1,15)=\gcd(5,15)=5.
\]
We recovered $15 = 3\times 5$ purely from the period $r=4$. A quantum computer does not do the period-finding by computing the powers one at a time the way we just did --- that would be no faster than classical. Instead, it uses interference.
\end{example}

\subsection{Finding the period with a quantum Fourier transform}
\label{ch35:qft}

The quantum part has a clean shape. Prepare a uniform superposition over a large range of inputs $x \in \{0,1,\dots,Q-1\}$ (with $Q=2^q$ a power of two, far bigger than $N$), and in a second register compute $a^x \bmod N$ \emph{reversibly}, entangling input and output:
\[
  \frac{1}{\sqrt{Q}}\sum_{x=0}^{Q-1}\ket{x}\,\ket{a^x \bmod N}.
\]
Measuring the second register collapses the first to a superposition over only those $x$ giving the observed value $a^{x_0}\bmod N$ --- and those $x$ are exactly an arithmetic progression $x_0, x_0+r, x_0+2r,\dots$ spaced by the period $r$. Now apply the \emph{quantum Fourier transform} to the first register. The QFT of a spike train spaced by $r$ is a spike train spaced by $\approx Q/r$: measuring yields a value close to a multiple of $Q/r$. From a measured $c \approx kQ/r$ you recover $r$ by rounding $c/Q$ to a nearby fraction with small denominator (continued fractions). Repeat a couple of times to nail down $k$ and $r$.

\begin{whybox}
Quantum interference is what makes the period pop out in one shot. Say $Q = 16$ and the period is $r = 4$, so collapsing the second register leaves four spikes at $x \in \{1, 5, 9, 13\}$. At frequency $c = 4 = Q/r$, each step of $\Delta x = 4$ advances the QFT phase by $4 \cdot 4 / 16 = 1$ full turn, so all four amplitude arrows point in the exact same direction and add to $4\times$ the length ($16\times$ the probability). At an off-period frequency like $c = 1$, each step rotates by $4/16 = 1/4$ turn ($+1, +i, -1, -i$), and the four arrows cancel to zero. Classically, finding $r$ takes $\Theta(r)$ evaluations where $r$ can be near $N$; the quantum circuit needs only several thousand logical qubits and $O((\log N)^3)$ gates for a $2048$-bit $N$, so doubling RSA key sizes cannot outrun it.
\end{whybox}

\begin{figure}[H]
\centering
\begin{tikzpicture}
  % Panel A: Modular Exponentiation Sequence
  \begin{scope}[shift={(0,0)}]
    \node[font=\small\bfseries] at (2.7, 3.6) {Stage 1: Superposition \& Periodicity};
    
    % Waveform
    \draw[->, thick] (0,0) -- (5.4,0) node[right, font=\tiny] {$x$};
    \draw[->, thick] (0,0) -- (0,3.0) node[above, font=\tiny] {$a^x \bmod N$};
    
    % Periodic points
    \foreach \k in {0, 1, 2, 3} {
      \pgfmathsetmacro{\xzero}{\k*1.3 + 0.2}
      \pgfmathsetmacro{\xone}{\k*1.3 + 0.5}
      \pgfmathsetmacro{\xtwo}{\k*1.3 + 0.8}
      \pgfmathsetmacro{\xthree}{\k*1.3 + 1.1}
      \filldraw[sourceblue!80!black] (\xzero, 0.6) circle (1.5pt);
      \filldraw[sourceblue!80!black] (\xone, 2.2) circle (1.5pt);
      \filldraw[sourceblue!80!black] (\xtwo, 1.4) circle (1.5pt);
      \filldraw[sourceblue!80!black] (\xthree, 0.4) circle (1.5pt);
    }
    \draw[dashed, sourceblue!50] (0.2, 0.6) -- (1.5, 0.6) -- (2.8, 0.6) -- (4.1, 0.6);
    \draw[dashed, trapred, thick] (0.2, 0) -- (0.2, 2.6);
    \draw[dashed, trapred, thick] (1.5, 0) -- (1.5, 2.6);
    \draw[{Stealth[length=2mm]}-{Stealth[length=2mm]}, trapred, thick] (0.2, 2.5) -- node[above, font=\tiny\bfseries] {Period $r$} (1.5, 2.5);
    \node[font=\tiny, fill=white, inner sep=1pt] at (2.7, -0.4) {Hidden period: $a^{x+r} \equiv a^x \pmod N$};
  \end{scope}

  % Panel B: QFT Interference Peaks
  \begin{scope}[shift={(7.2, 0)}]
    \node[font=\small\bfseries] at (2.7, 3.6) {Stage 2: QFT Interference Peaks};
    
    \draw[->, thick] (0,0) -- (5.4,0) node[right, font=\tiny] {Measured $c$};
    \draw[->, thick] (0,0) -- (0,3.0) node[above, font=\tiny] {Probability};
    
    % Peaks at multiples of Q/r
    \foreach \p in {0.5, 1.8, 3.1, 4.4} {
      \draw[structurepurple, very thick] (\p, 0) -- (\p, 2.4);
      \filldraw[structurepurple] (\p, 2.4) circle (2pt);
    }
    \draw[{Stealth[length=2mm]}-{Stealth[length=2mm]}, structurepurple, thick] (1.8, 2.1) -- node[above, font=\tiny\bfseries] {$\approx Q/r$} (3.1, 2.1);
    \node[structurepurple, font=\tiny] at (0.5, -0.3) {$0$};
    \node[structurepurple, font=\tiny] at (1.8, -0.3) {$\frac{Q}{r}$};
    \node[structurepurple, font=\tiny] at (3.1, -0.3) {$\frac{2Q}{r}$};
    \node[structurepurple, font=\tiny] at (4.4, -0.3) {$\frac{3Q}{r}$};

    \node[draw=targetgreen!80, fill=targetgreen!12, rounded corners=2pt, font=\tiny, align=center, inner sep=3pt] at (2.6, -0.9) {
      \textbf{Factor Extraction:}\\[1pt]
      Continued fractions yields $r \implies \gcd(a^{r/2} \pm 1, N) \text{ reveals factors}$
    };
  \end{scope}
\end{tikzpicture}
\caption{Shor's quantum period-finding architecture. In Stage 1, the quantum register evaluates $a^x \bmod N$ across an exponential superposition, establishing a periodic spike train spaced by $r$. In Stage 2, the Quantum Fourier Transform causes destructive interference everywhere except at discrete frequency multiples of $Q/r$, allowing measurement to disclose the order $r$ and break RSA in polynomial time.}
\label{ch35:fig-shor-qft}
\end{figure}

\subsection{Discrete log dies the same way}
\label{ch35:shor-dlog}

The discrete-log problem --- given $g$ and $y=g^s$ in a cyclic group of order $n$, find $s$ --- is also a hidden-period problem. Define $f(x_1,x_2)=g^{x_1} y^{-x_2}$. This function is constant on the set of $(x_1,x_2)$ with $x_1 - s x_2 \equiv 0$, i.e.\ it has a hidden \emph{linear} period determined by $s$. The same Fourier machinery extracts $s$ in polynomial time. Because elliptic-curve discrete log lives in a group with the same abelian structure, ECC offers \emph{no} extra protection here: a 256-bit elliptic-curve key and a 3072-bit RSA key are both Shor-breakable, and the elliptic-curve case actually needs fewer logical qubits because the numbers are smaller. The hardness that made ECC keys short is exactly the hardness Shor dissolves.

\begin{focusbox}
The pattern: Shor breaks any cryptosystem whose security reduces to finding a \emph{hidden period} or \emph{hidden subgroup} in a \emph{commutative} (abelian) group. Factoring, RSA, finite-field discrete log, elliptic-curve discrete log --- all of these. The post-quantum schemes in this chapter are built so that their hardness does \emph{not} have this abelian-hidden-subgroup shape. That is the structural reason lattices, codes, and hash chains are candidates: no one knows how to phrase their hard problems as a period to be found.
\end{focusbox}

\section{Grover's algorithm and what it does to symmetric crypto}
\label{ch35:grover}

Where Shor exploits algebraic structure for an exponential speedup, Grover searches without structure and achieves only a quadratic one. Grover's algorithm searches an unstructured space of $M$ possibilities for a marked item in $\Theta(\sqrt{M})$ quantum queries, versus $\Theta(M)$ classically. Applied to a block cipher with a $k$-bit key, brute force has $M=2^k$ candidates, so Grover finds the key in about $2^{k/2}$ steps.

The mechanism, in one paragraph: start in the uniform superposition over all $M$ keys, amplitude $1/\sqrt M$ each. Repeat a two-step ``rotation'' --- (1) an oracle that flips the sign of the amplitude of the correct key, (2) a ``diffusion'' that reflects all amplitudes about their average. Geometrically, each repetition rotates the state a tiny angle $\approx 2/\sqrt M$ toward the target; after about $(\pi/4)\sqrt M$ rotations the state is almost entirely on the target and a measurement returns it. Too many rotations \emph{overshoot} and the success probability falls again --- another reminder that this is rotation, not accumulation.

\begin{whybox}
Why only a square-root speedup, not exponential? Because there is no structure to exploit --- the search is over a black box. Each Grover step rotates by an angle proportional to $1/\sqrt M$, and you need to cover a constant angle, so you need $\Theta(\sqrt M)$ steps. A theorem (the BBBV bound, 1997) proves this is optimal: no quantum algorithm searches an unstructured space of size $M$ in fewer than $\Omega(\sqrt M)$ queries. So $2^{k/2}$ is a hard floor for generic key search, not just the best we currently know.
\end{whybox}

The defensive consequence is mild and you already know the fix: \emph{double the key size}. A $128$-bit key offers $\approx 2^{64}$ Grover security, which is uncomfortable; a $256$-bit key offers $\approx 2^{128}$, which is fine. This is why AES-256, not AES-128, is the conservative post-quantum choice, and why SHA-512 / SHA-256 hashing is judged adequate: the best quantum collision and preimage attacks give at most quadratic (and for collisions, by some analyses, only cube-root) speedups, leaving large hashes comfortable.

\begin{trapbox}
A common over-reaction: ``Grover halves every security level, so all symmetric crypto is in trouble.'' Two corrections. First, halving the \emph{exponent} (from $2^{128}$ work to $2^{64}$) is dramatic for $128$-bit keys but readily fixed by going to $256$ bits --- symmetric keys are cheap. Second, Grover is deeply serial: each of the $2^{k/2}$ steps depends on the last, so unlike classical brute force you cannot parallelize it into a $2^{k/2}$ speedup across many machines. Splitting Grover across $p$ processors only gives a $\sqrt p$ improvement, so the real-world cost of a Grover attack on AES-256 is astronomically beyond any plausible machine. The asymmetry is the headline: Shor is catastrophic and unfixable for RSA/ECC; Grover is a nuisance, fixed by one parameter bump.
\end{trapbox}

\begin{connectionbox}
The square-root in Grover is the same square-root you met in the companion probability book as the birthday bound, but turned inside out. There, $\sqrt M$ random samples suffice to find a \emph{collision} among $M$ values because collisions are quadratically common. Here, $\sqrt M$ \emph{queries} suffice to find a single marked item because amplitude rotation accumulates as the square root of the number of steps. Both are statements that ``$\sqrt M$ is the natural scale of $M$ possibilities'' --- one probabilistic, one geometric.
\end{connectionbox}

\begin{table}[H]
\centering
\begin{tabular}{@{}lll@{}}
\toprule
Primitive & Classical security & Best quantum attack \\
\midrule
RSA-2048 / ECC-256 & sub-exponential / $2^{128}$ & polynomial (Shor) --- broken \\
AES-128 & $2^{128}$ & $\approx 2^{64}$ (Grover) --- weak \\
AES-256 & $2^{256}$ & $\approx 2^{128}$ (Grover) --- safe \\
SHA-256 (collision) & $2^{128}$ & $\approx 2^{85}$ or $2^{128}$, model-dependent --- safe \\
SHA-256 (preimage) & $2^{256}$ & $\approx 2^{128}$ (Grover) --- safe \\
\bottomrule
\end{tabular}
\caption{Quantum impact at a glance. Public-key is the emergency; symmetric is a parameter bump.}
\label{ch35:tab-impact}
\end{table}

\section{Lattices: the foundation of post-quantum public key}
\label{ch35:lattices}

If period-finding is the abelian structure that Shor exploits, we want hard problems with \emph{no} such structure. Lattices deliver, and they are the basis for most of what NIST standardized. Start with the geometry, then build the cryptography.

\subsection{What a lattice is}
\label{ch35:lattice-def}

Given $n$ linearly independent vectors $\mathbf{b}_1,\dots,\mathbf{b}_n$ in $\R^n$, the \emph{lattice} they generate is the set of all integer combinations:
\[
  \mathcal{L} = \set[\Big]{\textstyle\sum_{i=1}^n z_i \mathbf{b}_i : z_i \in \Z}.
\]
It is an infinite grid of points --- regular, but in general skewed. The same lattice has many bases: you can replace the $\mathbf{b}_i$ by other integer combinations and still hit exactly the same points. Some bases are ``nice'' (short, nearly orthogonal vectors); some are ``ugly'' (long, nearly parallel vectors). The gap between nice and ugly bases is where the cryptography lives.

Two hard problems:
\begin{itemize}
  \item \textbf{Shortest Vector Problem (SVP):} find the shortest nonzero vector in $\mathcal L$.
  \item \textbf{Closest Vector Problem (CVP):} given an arbitrary target point $\mathbf t \in \R^n$ (not necessarily in the lattice), find the lattice point nearest to it.
\end{itemize}
In two or three dimensions these look easy --- you can see the answer geometrically. In a few hundred dimensions, with an ugly skewed basis, they are computationally formidable: the best known classical algorithms require time exponential in $n$, and \emph{no quantum speedup beyond polynomial factors is known}. There is no periodic abelian group structure for Shor's quantum Fourier transform to exploit.

\begin{figure}[H]
\centering
\begin{tikzpicture}
  % Panel A: Good vs Bad Basis
  \begin{scope}[shift={(0,0)}]
    \node[font=\small\bfseries] at (1.8, 3.6) {Panel A: Lattice Bases (Good vs Bad)};
    
    % Grid points
    \foreach \i in {-1,0,1,2,3} {
      \foreach \j in {-1,0,1,2} {
        \pgfmathsetmacro{\x}{\i*1.1 + \j*0.45}
        \pgfmathsetmacro{\y}{\j*1.1 + \i*0.15}
        \fill[gray!50] (\x,\y) circle (1.5pt);
      }
    }
    
    % Good basis
    \draw[-{Stealth[length=2mm]}, very thick, sourceblue] (0,0) -- (1.1, 0.15) node[below right, font=\tiny\bfseries] {$\mathbf{b}_1$ (Short)};
    \draw[-{Stealth[length=2mm]}, very thick, sourceblue] (0,0) -- (0.45, 1.1) node[above left, font=\tiny\bfseries] {$\mathbf{b}_2$ (Orthogonal)};

    % Bad basis: c1 = 2b1 + b2 = (2.65, 1.4), c2 = 3b1 + 2b2 = (4.2, 2.65) or c1 = 2b1+b2=(2.65,1.4), c2 = 3b1+b2=(3.75,1.55)
    \draw[-{Stealth[length=2mm]}, very thick, trapred] (0,0) -- (2.65, 1.4) node[above left=1pt and -2pt, font=\tiny\bfseries, fill=white, inner sep=1pt] {$\mathbf{c}_1$ (Skewed)};
    \draw[-{Stealth[length=2mm]}, very thick, trapred] (0,0) -- (3.1, 0.7) node[below right=-1pt and 1pt, font=\tiny\bfseries, fill=white, inner sep=1pt] {$\mathbf{c}_2$ (Skewed)};

    % Target and closest point (at i=1, j=2 -> (2.0, 2.35))
    \filldraw[targetgreen!80!black] (1.6, 2.7) circle (2.2pt) node[above left, font=\tiny\bfseries] {Target $\mathbf{t}$};
    \draw[dashed, targetgreen!80!black, thick] (1.6, 2.7) -- (2.0, 2.35);
    \filldraw[structurepurple] (2.0, 2.35) circle (2.2pt) node[right=2pt, font=\tiny\bfseries] {Nearest $\mathbf{v}$};
  \end{scope}

  % Panel B: LWE Perturbation & Decoding Basin
  \begin{scope}[shift={(6.8, 0)}]
    \node[font=\small\bfseries] at (2.4, 3.6) {Panel B: LWE Error Perturbation};

    % Voronoi cell
    \draw[fill=sourceblue!8, draw=sourceblue!40, dashed, thick] (2.4, 1.6) circle (1.55cm);
    \node[sourceblue!80!black, font=\tiny] at (2.4, 2.75) {Safe Decoding Basin ($\|\mathbf{e}\| < q/4$)};

    % Exact point
    \filldraw[sourceblue!80!black] (2.4, 1.6) circle (2.5pt) node[below left, font=\scriptsize\bfseries] {$A\mathbf{s} \pmod q$};
    
    % Error vector
    \draw[-{Stealth[length=2.2mm]}, very thick, trapred] (2.4, 1.6) -- (3.35, 2.25) node[midway, above left, font=\tiny\bfseries, trapred] {Noise $\mathbf{e}$};
    
    % Noisy sample
    \filldraw[quotientorange!90!black] (3.35, 2.25) circle (2.5pt) node[above right, font=\scriptsize\bfseries] {$\mathbf{b} = A\mathbf{s} + \mathbf{e}$};

    \node[draw=targetgreen!80, fill=targetgreen!12, rounded corners=2pt, font=\tiny, align=center, inner sep=3pt] at (2.4, -0.5) {
      \textbf{Rounding Rule:}\\[1pt]
      If noise stays inside basin, rounding restores $A\mathbf{s}$.\\[1pt]
      Gaussian elimination explodes $\mathbf{e}$ across cells.
    };
  \end{scope}
\end{tikzpicture}
\caption{The geometric foundation of post-quantum lattice cryptography. Panel A contrasts a short, nearly orthogonal basis (secret key) with a long, skewed basis (public key), illustrating how finding the nearest lattice point becomes intractable with skewed axes. Panel B depicts the Learning With Errors (LWE) setting: a secret lattice coordinate is perturbed by a bounded noise vector $\mathbf{e}$ that conceals $\mathbf{s}$ from linear solvers while permitting clean threshold rounding during decryption.}
\label{ch35:fig-lattice-lwe}
\end{figure}

\begin{intuitionbox}
A good (short, orthogonal) basis is like a coordinate grid drawn on graph paper: to find the nearest lattice point to a target you just round each coordinate to the nearest integer. A bad (long, skew) basis describes the \emph{same} grid but with axes that are nearly parallel and very long; rounding in those coordinates lands you far from the true nearest point, because a small rounding error along a long, slanted axis throws you across many cells. The public key in a lattice scheme is the bad basis; the private key is the good basis. Anyone can describe the lattice; only the key-holder can solve CVP in it.
\end{intuitionbox}

\subsection{Learning With Errors (LWE), the assumption that does the work}
\label{ch35:lwe}

Modern lattice crypto rarely talks about bases directly. It uses an equivalent, friendlier formulation: \emph{Learning With Errors}, introduced by Oded Regev in 2005. Fix a modulus $q$ and a dimension $n$. The secret is a vector $\mathbf s \in \Zn{q}^n$. The adversary is given many noisy linear equations in $\mathbf s$:
\[
  \mathbf a_i \getsr \Zn{q}^n, \qquad b_i = \langle \mathbf a_i, \mathbf s\rangle + e_i \bmod q,
\]
where each $e_i$ is a small ``error'' drawn from a narrow distribution (say, $e_i \in \{-2,\dots,2\}$). The \emph{search-LWE} problem is to recover $\mathbf s$ from the pairs $(\mathbf a_i, b_i)$. The \emph{decision-LWE} problem is to distinguish such pairs from uniformly random pairs.

Without the errors this is trivial: $n$ exact equations $b_i=\langle \mathbf a_i,\mathbf s\rangle$ are a linear system, and Gaussian elimination recovers $\mathbf s$. The errors are everything. Gaussian elimination amplifies the small errors catastrophically --- each elimination step adds errors together, and after $n$ steps the noise swamps the signal mod $q$. Regev proved that solving LWE on \emph{average} is as hard as solving lattice problems (like approximate SVP) in the \emph{worst case} --- a strong theoretical guarantee that ties the security of a randomly generated key to the hardness of the hardest instance of a geometric problem.

\begin{example}[A tiny LWE instance]
\label{ch35:ex-lwe}
Let $q=17$, $n=3$, secret $\mathbf s=(4,1,11)$, errors in $\{-1,0,1\}$. Three samples:
\[
\begin{aligned}
\mathbf a_1=(1,0,2): &\quad \langle\mathbf a_1,\mathbf s\rangle = 4+0+22=26\equiv 9,\ e_1=1 \Rightarrow b_1=10.\\
\mathbf a_2=(3,2,1): &\quad \langle\mathbf a_2,\mathbf s\rangle = 12+2+11=25\equiv 8,\ e_2=-1 \Rightarrow b_2=7.\\
\mathbf a_3=(0,5,4): &\quad \langle\mathbf a_3,\mathbf s\rangle = 0+5+44=49\equiv 15,\ e_3=0 \Rightarrow b_3=15.
\end{aligned}
\]
Handed only the rows $\mathbf a_i$ and the values $b_i$, you must find $\mathbf s$. With $q=17$ and $n=3$ you could brute-force all $17^3=4913$ candidates and pick the one whose errors are all small --- that is the attack, and it is exactly why real parameters use $n\approx 256$ or more so that $q^n$ is astronomically large. The brute force you \emph{can} do here is the brute force that becomes infeasible at scale.
\end{example}

\subsection{Encrypting with LWE: Regev's scheme}
\label{ch35:regev}

LWE gives encryption almost for free. Stack $m$ samples into a matrix $A\in \Zn{q}^{m\times n}$ and a vector $\mathbf b = A\mathbf s + \mathbf e \bmod q$. The public key is $(A,\mathbf b)$; the secret key is $\mathbf s$. To encrypt one bit $\mu\in\{0,1\}$:

\begin{enumerate}
  \item Pick a random subset of the rows --- equivalently a random vector $\mathbf r\in\{0,1\}^m$.
  \item Send the ciphertext $\big(\mathbf u, v\big) = \big(A^{\!\top}\mathbf r,\ \langle \mathbf b,\mathbf r\rangle + \mu\cdot\floor{q/2}\big) \bmod q$.
\end{enumerate}
To decrypt, compute $v - \langle \mathbf s, \mathbf u\rangle \bmod q$. Expand:
\[
  v - \langle\mathbf s,\mathbf u\rangle = \langle A\mathbf s+\mathbf e,\mathbf r\rangle + \mu\floor{q/2} - \langle \mathbf s, A^{\!\top}\mathbf r\rangle = \langle\mathbf e,\mathbf r\rangle + \mu\floor{q/2}.
\]
The term $\langle\mathbf e,\mathbf r\rangle$ is a sum of a few small errors, so it is small relative to $q/2$. The result is therefore close to $0$ (if $\mu=0$) or close to $q/2$ (if $\mu=1$); rounding recovers $\mu$. On Example~\ref{ch35:ex-lwe}'s numbers ($q = 17$, $\mathbf{s} = (4,1,11)$), encrypting $\mu = 1$ ($\floor{17/2} = 8$) with subset $\mathbf{r} = (1, 0, 1)$ sums rows $1$ and $3$: $\mathbf{u} = (1,0,2) + (0,5,4) = (1,5,6)$ and $v = 10 + 15 + 8 = 33 \equiv 16 \pmod{17}$. Decrypting computes $v - \langle\mathbf{s}, \mathbf{u}\rangle = 16 - (4 + 5 + 66) \equiv 16 - 7 = 9 \pmod{17}$ --- which is $e_1 + e_3 + 8 = 1 + 0 + 8 = 9$, sitting right next to $8$ and rounding cleanly to $\mu = 1$. Figure~\ref{ch35:fig-regev-flow} traces this flow from public-key generation through subset-sum encryption and noisy cancellation. Security: distinguishing the ciphertext from random is exactly decision-LWE, so an eavesdropper who could read $\mu$ could solve LWE.

\begin{figure}[H]
\centering
\begin{tikzpicture}[font=\small,
  actor/.style={draw=sourceblue!80, fill=sourceblue!10, rounded corners=3pt, minimum width=3.4cm, minimum height=0.8cm, align=center},
  stepbox/.style={draw=gray!60, fill=gray!6, rounded corners=3pt, minimum width=3.6cm, inner sep=5pt, align=center, font=\footnotesize},
  arr/.style={-{Stealth[length=2.4mm]}, thick, sourceblue!80!black}
]
  % Columns: Alice (Receiver/KeyGen & Decrypt) and Bob (Sender/Encrypt)
  \node[actor, draw=aliceblue!80, fill=aliceblue!12] (alice) at (0, 3.6) {\textbf{Alice} (Key Holder)};
  \node[actor, draw=bobgreen!80, fill=bobgreen!12] (bob) at (7.4, 3.6) {\textbf{Bob} (Sender, bit $\mu\in\{0,1\}$)};

  % Step 1: KeyGen
  \node[stepbox, draw=aliceblue!60] (kg) at (0, 2.2) {
    \textbf{1. Key Generation}\\[2pt]
    Secret $\mathbf{s}\in\Zn{q}^n$, small noise $\mathbf{e}$\\
    $\mathbf{b} = A\mathbf{s} + \mathbf{e} \bmod q$
  };

  % Step 2: Encrypt
  \node[stepbox, draw=bobgreen!60] (enc) at (7.4, 2.2) {
    \textbf{2. Subset-Sum Encryption}\\[2pt]
    Pick random subset $\mathbf{r}\in\{0,1\}^m$\\
    $\mathbf{u} = A^{\!\top}\mathbf{r}$, \quad $v = \langle\mathbf{b},\mathbf{r}\rangle + \mu\floor{q/2}$
  };

  \draw[arr, aliceblue!80!black] (kg.east) -- node[above, font=\scriptsize\bfseries] {Public key $(A,\mathbf{b})$} (enc.west);

  % Step 3: Decrypt
  \node[stepbox, draw=targetgreen!80, fill=targetgreen!8, minimum width=4.4cm] (dec) at (0, 0.2) {
    \textbf{3. Noisy Cancellation \& Rounding}\\[2pt]
    $v - \langle\mathbf{s},\mathbf{u}\rangle = \underbrace{\langle A\mathbf{s},\mathbf{r}\rangle - \langle\mathbf{s},A^{\!\top}\mathbf{r}\rangle}_{=0} + \underbrace{\langle\mathbf{e},\mathbf{r}\rangle}_{\text{small}} + \mu\floor{q/2}$\\[2pt]
    Near $0 \Rightarrow \mu=0$; \quad near $q/2 \Rightarrow \mu=1$
  };

  \draw[arr, bobgreen!80!black] (enc.south) |- node[pos=0.72, above, font=\scriptsize\bfseries] {Ciphertext $(\mathbf{u},v)$} (dec.east);
\end{tikzpicture}
\caption{Regev's LWE encryption and decryption pipeline. Bob masks his message bit $\mu\floor{q/2}$ with a random subset sum of Alice's noisy public equations. When Alice subtracts $\langle\mathbf{s},\mathbf{u}\rangle$, the large matrix terms $A\mathbf{s}$ cancel identically, leaving only the small accumulated noise $\langle\mathbf{e},\mathbf{r}\rangle$ on top of the encoded bit.}
\label{ch35:fig-regev-flow}
\end{figure}

\begin{whybox}
The error is doing two opposite jobs at once, and the design is the balance between them. It must be \emph{large enough} that the equations are hard --- so an attacker cannot solve the linear system --- yet \emph{small enough} that the legitimate decryptor's noise term $\langle\mathbf e,\mathbf r\rangle$ stays well under $q/4$, leaving room to round correctly. Every lattice scheme is a budget: noise growth (which protects you) versus correctness margin (which serves you). Push the error up for security and you risk decryption failures; push it down for reliability and you weaken the assumption. Choosing $q$, $n$, and the error width is exactly negotiating that budget. Decryption-failure probability even becomes an attack surface --- an adversary who can provoke and observe failures learns about $\mathbf s$ --- which is why standardized parameters drive the failure rate down to around $2^{-130}$ or lower.
\end{whybox}

\subsection{Module-LWE, Kyber, and Dilithium}
\label{ch35:kyber-dilithium}

Plain LWE has a practical problem: the public key is a full $m\times n$ matrix of integers mod $q$, which is large (megabytes). The fix is to add a controlled dose of structure, trading a little theoretical conservatism for a huge size win. Replace integers mod $q$ with elements of a polynomial ring $R_q = \Zn{q}[X]/(X^N+1)$, where $X^N \equiv -1$ --- a single element is now a degree-$(N-1)$ polynomial, and multiplying by $X$ rotates all $N$ coefficients by one slot while negating the wrapped term (this is \emph{Ring-LWE}). One ring element does the work of an $N\times N$ structured matrix, so keys shrink dramatically. \emph{Module-LWE} sits in between: small matrices ($2\times2$, $3\times3$, $4\times4$) of ring elements, giving a tunable knob between size and conservatism. This is what the NIST winners use.

\begin{itemize}
  \item \textbf{ML-KEM (Kyber)} --- FIPS 203, standardized 2024 --- is a \emph{key-encapsulation mechanism}: it lets two parties agree on a shared secret, the post-quantum replacement for Diffie-Hellman / ECDH. It is Module-LWE encryption (essentially the Regev scheme above, over a module) wrapped with the Fujisaki-Okamoto transform to upgrade from ``secure against eavesdroppers'' to ``secure against active, chosen-ciphertext attackers.'' Kyber-768 targets the AES-192 security level with a public key of about $1.2$\,KB and a ciphertext of about $1.1$\,KB --- small enough to drop into a TLS handshake.
  \item \textbf{ML-DSA (Dilithium)} --- FIPS 204, standardized 2024 --- is a \emph{signature} scheme over Module-LWE, the post-quantum replacement for ECDSA/RSA signatures. It uses Vadim Lyubashevsky's ``Fiat-Shamir with aborts'' technique. Why the aborts? In classical Schnorr over $\Z_q$, a uniform mask $y$ hides $s$ in $z = y + cs \bmod q$ because $\Z_q$ wraps cyclically with no edges. In a lattice, the response $\mathbf{z} = \mathbf{y} + c\mathbf{s}$ must stay \emph{short} without mod-$q$ wraparound: if $\mathbf{y}$ is drawn from a bounded box $[-B, B]^n$, adding $c\mathbf{s}$ shifts that box, and points near the boundary leak $\mathbf{s}$. The signer therefore \emph{rejects and retries} with a fresh $\mathbf{y}$ whenever any coordinate of $\mathbf{z}$ lands within $\beta$ of the box edge ($\|\mathbf{z}\|_\infty \ge B - \beta$), cropping away the shifted fringes so every published $\mathbf{z}$ is uniform on the safe inner box $[-(B-\beta), B-\beta]^n$ and independent of $\mathbf{s}$. Signatures are a couple of kilobytes --- larger than ECDSA's $64$ bytes, the recurring cost of going post-quantum.
\end{itemize}

\begin{focusbox}
Because NIST renamed the competition submissions upon standardization, you will meet both sets of names in libraries and papers: \textbf{Kyber} is \textbf{ML-KEM} (FIPS~203), \textbf{Dilithium} is \textbf{ML-DSA} (FIPS~204), \textbf{SPHINCS+} is \textbf{SLH-DSA} (FIPS~205), and \textbf{Falcon} is \textbf{FN-DSA} (draft FIPS~206). Module-lattice schemes won the primary slots because they hit the best balance of key size, speed, and confidence. Still, structured lattices assume that Module-LWE is as hard as unstructured LWE --- an assumption with weaker worst-case backing --- which is why NIST also standardized hash-based SLH-DSA and continued evaluating code-based KEMs as a hedge.
\end{focusbox}

\begin{attackbox}
The danger of \emph{too much} structure is not hypothetical. In 2022 Castryck and Decru broke SIKE --- a leading isogeny-based KEM that had reached round 4 of the NIST process --- in about an hour on a single classical core, using a theorem (Kani's) about products of elliptic curves that exploited auxiliary structure SIKE exposed. SIKE was not a lattice scheme, but the lesson generalizes: each bit of algebraic structure you add to shrink keys is a bit an attacker might find a handle on. This is the live tension in Ring-/Module-LWE. No subexponential attack on them is known, but the extra structure means the worst-case-to-average-case guarantee is weaker than for unstructured LWE, and prudent designs (Kyber among them) keep the structure modest --- small modules rather than a single huge ring --- precisely to stay away from the SIKE cliff.
\end{attackbox}

\section{A toy lattice cryptosystem you can break by hand}
\label{ch35:toy-lattice}

To feel where the security lives, here is a deliberately tiny LWE encryption and the brute-force attack on it. \textbf{This is for learning only --- the parameters are minuscule and instantly breakable, and you must never use anything like it to protect real data.}

\begin{implementbox}
\begin{lstlisting}[language=Python]
# TOY LWE encryption of a single bit. LEARNING ONLY -- INSECURE.
# Real schemes (ML-KEM/Kyber) use n>=256, structured rings,
# the Fujisaki-Okamoto transform, and constant-time code.
import random, itertools

q = 97          # modulus (tiny on purpose)
n = 6           # secret dimension (tiny on purpose)
m = 12          # number of LWE samples in the public key

def dot(u, v):  return sum(a*b for a, b in zip(u, v)) % q
def small():    return random.randint(-1, 1)   # narrow error

def keygen():
    s = [random.randrange(q) for _ in range(n)]            # secret
    A = [[random.randrange(q) for _ in range(n)] for _ in range(m)]
    b = [(dot(A[i], s) + small()) % q for i in range(m)]   # noisy
    return (A, b), s                                       # pk, sk

def encrypt(pk, bit):
    A, b = pk
    r = [random.randint(0, 1) for _ in range(m)]           # subset
    u = [sum(r[i]*A[i][j] for i in range(m)) % q for j in range(n)]
    v = (sum(r[i]*b[i] for i in range(m)) + bit*(q//2)) % q
    return u, v

def decrypt(sk, ct):
    u, v = ct
    w = (v - dot(sk, u)) % q          # = <e,r> + bit*(q//2)
    if w > q//2: w -= q              # fold into [-q/2, q/2)
    return 1 if abs(w) > q//4 else 0

pk, sk = keygen()
for bit in (0, 1):
    assert decrypt(sk, encrypt(pk, bit)) == bit
print("toy LWE round-trips correctly")

# ---- the attack: brute force the secret because n,q are tiny ----
def attack(pk):
    A, b = pk
    for s in itertools.product(range(q), repeat=n):   # q**n guesses
        if all(abs(((dot(A[i], s) - b[i] + q//2) % q) - q//2) <= 1
               for i in range(m)):
            return s    # a secret consistent with small noise
# attack(pk) recovers sk because q**n (here 97^6) is small enough.
# At n=256 this loop has q**256 iterations -- utterly infeasible.
\end{lstlisting}
\end{implementbox}

The attack is the lesson. With $n=6$ and $q=97$ there are $97^6 \approx 8.3\times 10^{11}$ candidate secrets --- a few minutes of CPU, or far less with smarter lattice reduction. Every defense in real lattice crypto is about pushing that number out of reach \emph{while keeping the legitimate error budget intact}: bigger $n$, careful error distributions, and structure (rings/modules) so the keys stay small even as $n$ climbs into the hundreds.

\section{Hash-based signatures: security from the primitive you trust most}
\label{ch35:hash-sigs}

Lattices are efficient but young; their hardness assumptions are a couple of decades old. There is one corner of post-quantum cryptography that asks you to trust \emph{nothing new}: hash-based signatures. Their only assumption is that a hash function is a good one-way / collision-resistant function --- and Grover merely halves that, so a $256$-bit hash is fine. If you believe in SHA-256 (or SHA-3), you believe in these signatures, full stop. The cost is awkwardness: state, size, and a limited number of signatures per key.

\subsection{Lamport: signing one bit, then one message, with a hash}
\label{ch35:lamport}

Leslie Lamport's 1979 one-time signature is the seed of the whole family. To sign a single bit:
\begin{itemize}
  \item \textbf{Secret key:} two random values $x_0, x_1$.
  \item \textbf{Public key:} their hashes $y_0=\Hash(x_0),\ y_1=\Hash(x_1)$, published.
  \item \textbf{Sign bit $b$:} reveal $x_b$.
  \item \textbf{Verify:} check $\Hash(x_b)=y_b$.
\end{itemize}
Revealing the preimage of a hash that only you knew is the signature; preimage resistance is what stops a forger. To sign an $\ell$-bit message digest, use $\ell$ such pairs: publish $2\ell$ hashes, and for each message bit reveal the matching half of the corresponding pair.

\begin{example}[Lamport on 4 bits]
\label{ch35:ex-lamport}
Suppose the message digest is $1011$. The secret key is four pairs $(x_{i,0},x_{i,1})$ for $i=1,\dots,4$; the public key is the eight hashes $y_{i,b}=\Hash(x_{i,b})$. As Figure~\ref{ch35:fig-lamport} shows, the signature for $1011$ is $(x_{1,1},\,x_{2,0},\,x_{3,1},\,x_{4,1})$ --- pick the $1$-half, the $0$-half, the $1$-half, the $1$-half. The verifier hashes each revealed value and checks it equals the published hash for that bit position and bit value. A forger who wants to sign a \emph{different} digest must produce a preimage of some $y_{i,b}$ they were never shown --- exactly a preimage attack on the hash.
\end{example}

\begin{figure}[H]
\centering
\begin{tikzpicture}[font=\footnotesize,
  skbox/.style={draw=gray!60, fill=gray!8, rounded corners=2pt, minimum width=1.2cm, minimum height=0.55cm},
  revbox/.style={draw=targetgreen!80!black, fill=targetgreen!18, very thick, rounded corners=2pt, minimum width=1.2cm, minimum height=0.55cm, font=\footnotesize\bfseries},
  pkbox/.style={draw=sourceblue!70, fill=sourceblue!8, rounded corners=2pt, minimum width=1.2cm, minimum height=0.55cm},
  arr/.style={-{Stealth[length=1.8mm]}, thick, targetgreen!75!black}
]
  \node[anchor=east, font=\small\bfseries] at (-1.5, 2.3) {Message bits:};
  \node[anchor=east, font=\small\bfseries] at (-1.5, 1.3) {Secret key $sk$:};
  \node[anchor=east, font=\small\bfseries] at (-1.5, -0.3) {Public key $pk$:};

  % Bit 1: b=1
  \node[font=\small\bfseries, targetgreen!75!black] at (0.7, 2.3) {$b_1 = 1$};
  \node[skbox] (x10) at (0, 1.55) {$x_{1,0}$};
  \node[revbox] (x11) at (1.4, 1.55) {$x_{1,1}$};
  \node[pkbox] (y10) at (0, -0.3) {$y_{1,0}$};
  \node[pkbox, draw=targetgreen!80!black, very thick] (y11) at (1.4, -0.3) {$y_{1,1}$};
  \draw[-{Stealth[length=1.8mm]}, gray!60] (x10) -- (y10);
  \draw[arr] (x11) -- node[right, font=\tiny] {$\Hash$} (y11);

  % Bit 2: b=0
  \node[font=\small\bfseries, targetgreen!75!black] at (3.7, 2.3) {$b_2 = 0$};
  \node[revbox] (x20) at (3.0, 1.55) {$x_{2,0}$};
  \node[skbox] (x21) at (4.4, 1.55) {$x_{2,1}$};
  \node[pkbox, draw=targetgreen!80!black, very thick] (y20) at (3.0, -0.3) {$y_{2,0}$};
  \node[pkbox] (y21) at (4.4, -0.3) {$y_{2,1}$};
  \draw[arr] (x20) -- node[right, font=\tiny] {$\Hash$} (y20);
  \draw[-{Stealth[length=1.8mm]}, gray!60] (x21) -- (y21);

  % Bit 3: b=1
  \node[font=\small\bfseries, targetgreen!75!black] at (6.7, 2.3) {$b_3 = 1$};
  \node[skbox] (x30) at (6.0, 1.55) {$x_{3,0}$};
  \node[revbox] (x31) at (7.4, 1.55) {$x_{3,1}$};
  \node[pkbox] (y30) at (6.0, -0.3) {$y_{3,0}$};
  \node[pkbox, draw=targetgreen!80!black, very thick] (y31) at (7.4, -0.3) {$y_{3,1}$};
  \draw[-{Stealth[length=1.8mm]}, gray!60] (x30) -- (y30);
  \draw[arr] (x31) -- node[right, font=\tiny] {$\Hash$} (y31);

  % Bit 4: b=1
  \node[font=\small\bfseries, targetgreen!75!black] at (9.7, 2.3) {$b_4 = 1$};
  \node[skbox] (x40) at (9.0, 1.55) {$x_{4,0}$};
  \node[revbox] (x41) at (10.4, 1.55) {$x_{4,1}$};
  \node[pkbox] (y40) at (9.0, -0.3) {$y_{4,0}$};
  \node[pkbox, draw=targetgreen!80!black, very thick] (y41) at (10.4, -0.3) {$y_{4,1}$};
  \draw[-{Stealth[length=1.8mm]}, gray!60] (x40) -- (y40);
  \draw[arr] (x41) -- node[right, font=\tiny] {$\Hash$} (y41);
\end{tikzpicture}
\caption{Lamport one-time signature on the 4-bit digest $1011$ (Example~\ref{ch35:ex-lamport}). For each bit $b_i$, the signer reveals the highlighted secret preimage $x_{i,b_i}$ while keeping the other half secret; the verifier hashes each revealed value and checks it against the corresponding public hash $y_{i,b_i}$.}
\label{ch35:fig-lamport}
\end{figure}

\begin{trapbox}
``One-time'' is not a slogan, it is a hard rule. Sign two different messages with one Lamport key and you leak both halves of some pair: where the two digests differ in bit $i$, you have revealed both $x_{i,0}$ and $x_{i,1}$. An attacker can then mix-and-match revealed halves to forge signatures on digests that combine bits from the two you signed. Reusing a one-time key is the hash-signature analogue of reusing a one-time-pad key or an ECDSA nonce: a catastrophic, total break, not a gradual erosion. This is why the entire engineering effort in this family is about safely getting \emph{many} signatures out of \emph{many} one-time keys.
\end{trapbox}

\subsection{From one signature to a tree: Merkle, XMSS, SPHINCS+}
\label{ch35:merkle-sphincs}

One Lamport key signs once. To sign $2^h$ messages, generate $2^h$ one-time key pairs and compress all their public keys into a single short public key using a \emph{Merkle tree}: hash the public keys into the leaves, hash pairs upward, and publish only the root. A signature now carries (a) a one-time signature from one leaf, plus (b) the \emph{authentication path} --- the sibling hashes along the route from that leaf to the root --- so the verifier can recompute the root and confirm the leaf belongs to your tree.

\begin{figure}[H]
\centering
\begin{tikzpicture}[
  treenode/.style={draw=structurepurple!80, fill=structurepurple!12, rounded corners=2pt, minimum width=2.0cm, minimum height=0.7cm, font=\small},
  pathnode/.style={draw=quotientorange!80, fill=quotientorange!18, rounded corners=2pt, minimum width=1.8cm, minimum height=0.65cm, font=\small\bfseries},
  activeleaf/.style={draw=targetgreen!80, fill=targetgreen!18, rounded corners=2pt, minimum width=1.5cm, minimum height=0.6cm, font=\footnotesize\bfseries},
  leaf/.style={draw=sourceblue!50, fill=sourceblue!8, rounded corners=2pt, minimum width=1.4cm, minimum height=0.55cm, font=\footnotesize},
  ln/.style={thick, gray!70},
  pathln/.style={very thick, quotientorange}
]
  \node[treenode, draw=sourceblue!80, fill=sourceblue!15] (root) at (0,3.2) {\textbf{Root Key} $pk = h_{0123}$};
  
  \node[treenode] (h01) at (-2.6, 1.8) {$h_{01}$};
  \node[pathnode] (h23) at (2.6, 1.8) {$h_{23}$ (Sibling)};
  
  \node[pathnode] (l0) at (-3.9, 0.4) {$pk_0$ (Sibling)};
  \node[activeleaf] (l1) at (-1.3, 0.4) {$pk_1$ (Signer)};
  \node[leaf] (l2) at (1.3, 0.4) {$pk_2$};
  \node[leaf] (l3) at (3.9, 0.4) {$pk_3$};
  
  \draw[pathln] (root) -- (h01);
  \draw[pathln] (root) -- (h23);
  \draw[pathln] (h01) -- (l0);
  \draw[pathln] (h01) -- (l1);
  \draw[ln] (h23) -- (l2);
  \draw[ln] (h23) -- (l3);

  \node[draw=quotientorange!80, fill=quotientorange!8, rounded corners=2pt, font=\scriptsize, align=center, inner sep=4pt] at (0, -0.7) {
    \textbf{Authentication Path for Leaf $pk_1$:}\\
    Signature contains one-time signature under $sk_1$ plus sibling path $\{pk_0, h_{23}\}$.\\
    Verifier computes $h_{01} = \Hash(pk_0 \mathbin{\Vert} pk_1)$, then verifies $\Hash(h_{01} \mathbin{\Vert} h_{23}) = pk$.
  };
\end{tikzpicture}
\caption{Merkle tree signature hierarchy. To authenticate message signatures across $2^h$ one-time Lamport key pairs, public keys are formed into a Merkle hash tree. When leaf $pk_1$ signs a message, the signer provides the one-time signature plus the authentication path ($pk_0$ and $h_{23}$). The verifier re-computes hashes up to the root, validating authenticity without trusting any public key except the single published root.}
\label{ch35:fig-merkle-sigs}
\end{figure}

The remaining wrinkle is \emph{state}: you must never reuse a leaf, so the signer has to remember which leaves are spent. A signer that forgets --- restored from a backup, cloned to two servers --- can reuse a leaf and break itself. \textbf{XMSS} and \textbf{LMS} (RFC 8391, RFC 8554) are the standardized stateful tree schemes; they are excellent where you can guarantee state, such as signing firmware images from a controlled build server. \textbf{SPHINCS+} --- standardized as \textbf{SLH-DSA}, FIPS 205 --- removes the state by picking the one-time leaf \emph{pseudorandomly from the message} and stacking trees of trees (a hypertree) so the chance of a damaging collision is negligible without bookkeeping. The price is size: SLH-DSA signatures run tens of kilobytes. You pay in bytes for the privilege of trusting only a hash.

\begin{applicationbox}
Stateful hash-based signatures are already the recommended tool for \emph{firmware and software signing}, where the threat is precisely ``harvest now, forge later'' over device lifetimes of decades, and where a controlled signing server can manage state safely. NIST's SP 800-208 approves XMSS and LMS for exactly this. SLH-DSA is the conservative fallback for general signatures when you want maximal confidence and can tolerate large signatures --- a belt-and-suspenders choice for long-lived roots of trust, such as the key that signs your other keys.
\end{applicationbox}

\section{Migration, hybrids, and timelines}
\label{ch35:migration}

\subsection{The standardized algorithms}
\label{ch35:nist}

As of 2024, the headline NIST standards are: \textbf{FIPS 203 (ML-KEM / Kyber)} for key establishment, \textbf{FIPS 204 (ML-DSA / Dilithium)} and \textbf{FIPS 205 (SLH-DSA / SPHINCS+)} for signatures, with \textbf{FN-DSA (Falcon)} following. NIST also ran a separate ``on-ramp'' for additional signatures to diversify away from lattices, and selected code-based and other KEM candidates (such as the conservative, large-key Classic McEliece) for niche use where key size matters less than decades of cryptanalytic confidence. The portfolio is deliberately not all-lattice: if a structural break hits lattices, the world wants a non-lattice scheme already on the shelf.

\subsection{Hybrid key exchange: belt and suspenders}
\label{ch35:hybrid}

Deploying a brand-new assumption (Module-LWE) as your \emph{only} protection is itself a risk --- young schemes have bugs and parameter mistakes, and a structural break is conceivable. The standard mitigation is a \emph{hybrid}: run a classical key exchange (X25519, say) \emph{and} a post-quantum one (ML-KEM) and combine both shared secrets through a key-derivation function:
\[
  K = \mathrm{KDF}\big(\,\text{X25519 secret} \concat \text{ML-KEM secret}\,\big).
\]
The session key is safe if \emph{either} ingredient is safe. A quantum computer breaks X25519 but not ML-KEM; a future cryptanalytic break of ML-KEM still leaves X25519, which holds against every classical adversary. You only lose if both fall, which is the conjunction of two independent failures (Figure~\ref{ch35:fig-hybrid-kdf}). This is exactly what shipped: TLS 1.3 deployments rolled out the hybrid \texttt{X25519MLKEM768} group across major browsers and servers starting in 2024, so a large fraction of web traffic already negotiates a post-quantum-protected handshake today.

\begin{figure}[H]
\centering
\begin{tikzpicture}[font=\small,
  box/.style={draw=sourceblue!80, fill=sourceblue!8, rounded corners=3pt, minimum width=3.8cm, minimum height=0.9cm, align=center},
  kdf/.style={draw=structurepurple!80, fill=structurepurple!12, rounded corners=3pt, minimum width=2.8cm, minimum height=1.2cm, align=center, font=\small\bfseries},
  outbox/.style={draw=targetgreen!80!black, fill=targetgreen!14, rounded corners=3pt, minimum width=3.2cm, minimum height=0.9cm, align=center, font=\small\bfseries},
  arr/.style={-{Stealth[length=2.2mm]}, thick}
]
  \node[box] (clas) at (0, 0.85) {\textbf{Classical KEM / DH}\\ \footnotesize e.g.\ X25519};
  \node[box, draw=quotientorange!85!black, fill=quotientorange!10] (pq) at (0, -0.85) {\textbf{Post-Quantum KEM}\\ \footnotesize e.g.\ ML-KEM-768};

  \node[kdf] (kdf) at (5.6, 0) {Combiner\\$\mathrm{KDF}(s_1 \concat s_2)$};
  \node[outbox] (sk) at (10.2, 0) {Session Key $K$\\ \footnotesize\normalfont Safe if \emph{either} holds};

  \draw[arr, sourceblue!80!black] (clas.east) -- node[above, font=\footnotesize] {$s_1$} (kdf.160);
  \draw[arr, quotientorange!85!black] (pq.east) -- node[below, font=\footnotesize] {$s_2$} (kdf.200);
  \draw[arr, targetgreen!80!black] (kdf.east) -- (sk.west);
\end{tikzpicture}
\caption{Hybrid key exchange (\texttt{X25519MLKEM768}). Classical and post-quantum shared secrets $s_1$ and $s_2$ are fed together into a key-derivation function, so the resulting session key $K$ remains indistinguishable from random as long as at least one of the two underlying problems stays hard.}
\label{ch35:fig-hybrid-kdf}
\end{figure}

\begin{applicationbox}
The hybrid is live in places you use daily. Chrome, Firefox, and Cloudflare enabled \texttt{X25519MLKEM768} (and its predecessor \texttt{X25519Kyber768}) for TLS, so a sizable share of HTTPS connections are already quantum-resistant on the key-exchange side. Signal deployed PQXDH --- a post-quantum extension of its X3DH key agreement that mixes a Kyber/ML-KEM secret into the initial handshake --- so the \emph{forward} secret of a Signal session resists future quantum decryption even though the long-term identity signatures are still classical. SSH (OpenSSH) defaults to a hybrid \texttt{sntrup761x25519} key exchange. The pattern is identical everywhere: protect the ephemeral key agreement first, because that is what ``harvest now, decrypt later'' targets; migrate signatures and certificates more slowly, because forging a signature requires the quantum computer to exist \emph{now}, not later.
\end{applicationbox}

\begin{focusbox}
Notice the priority ordering, and why it is correct. \emph{Confidentiality} is urgent because of store-and-decrypt: traffic recorded today is exposed the day the machine arrives, so the key exchange must go post-quantum first. \emph{Authentication} (signatures, certificates) is less urgent because a forged signature is only useful at the moment of the connection --- you cannot retroactively forge yesterday's handshake. An attacker cannot wait for a future quantum computer to impersonate a server in a session that already happened. So migrate KEMs now; migrate signature infrastructure (which is bigger and slower, given certificate chains and root programs) on a longer schedule.
\end{focusbox}

\subsection{Timelines, without hand-waving}
\label{ch35:timelines}

How long until a cryptographically relevant quantum computer (CRQC) exists? No one knows, and anyone who gives you a confident year is selling something. What we can say with precision is the engineering gap. Today's quantum processors have on the order of hundreds to low thousands of noisy physical qubits with error rates around $10^{-3}$. Breaking RSA-2048 with Shor needs a few thousand \emph{logical} qubits, and error correction inflates each logical qubit to roughly $10^3$--$10^4$ physical qubits, putting the requirement near \emph{millions} of high-quality physical qubits running stably for hours. That is many orders of magnitude beyond current hardware, and the gap is dominated by error rates, not raw qubit counts --- which is why ``qubit count'' headlines are a poor proxy for cryptographic danger.

Estimates from serious groups cluster in the range of one to three decades, with wide error bars and a real tail probability of ``sooner'' if error correction improves faster than expected. The defender's calculus does not actually depend on resolving that uncertainty. Mosca's inequality states it cleanly: if $X$ is how long your data must stay confidential, $Y$ is how long migration takes, and $Z$ is the time until a CRQC exists, then you have a problem whenever
\[
  X + Y > Z.
\]
For data that must stay secret twenty years, with a migration that realistically takes five to ten years across a large organization, $X+Y$ already reaches into the late 2040s --- comfortably inside the plausible range for $Z$. That is why migration is a present-day program, not a future one, even though the breaking machine is not here. The threat is real today not because the computer exists, but because the secrets and the recordings do.

\begin{trapbox}
Two failure modes bracket the sensible path. \emph{Complacency}: ``no quantum computer exists, so do nothing'' ignores store-and-decrypt and the multi-year reality of migration; by the time the machine is announced it is far too late to start. \emph{Panic}: ripping out classical crypto for unvetted post-quantum schemes deployed alone has already caused real incidents --- early non-standard parameter sets, implementation bugs, and at least one reference implementation with a timing side channel in the very rejection-sampling step that was supposed to add safety. The measured answer is the hybrid: deploy post-quantum \emph{alongside} classical, using standardized parameters and audited constant-time code, so a flaw in the new scheme cannot undo the protection of the old one while the new one matures.
\end{trapbox}

\section{Where this leaves us}
\label{ch35:summary}

Quantum computers threaten public-key cryptography in a specific, structural way: Shor's algorithm dissolves the abelian hidden-period problems --- factoring and discrete log --- that RSA and elliptic curves are built on, and no key size repairs that. Grover's algorithm only dents symmetric cryptography, and doubling the key size (AES-256, SHA-384/512) restores the margin. The replacements avoid abelian period structure: lattice problems (SVP/CVP, packaged as LWE and its ring/module variants) power the efficient new standards ML-KEM and ML-DSA, while hash-based signatures (XMSS, LMS, SLH-DSA) offer maximal confidence at the cost of size and, sometimes, state. The migration is already underway in TLS, Signal, and SSH, run as hybrids so that a stumble in the young schemes cannot cost you the protection of the old ones. The machine that forces all this is not here. The data it would read is.

With public-key primitives and their quantum-resistant successors in hand, Chapter~36 explores how cryptographic hashes, digital signatures, and consensus protocols combine to build decentralized state in blockchains.

\section{Exercises}
\label{ch35:exercises}

\begin{exercise}[Shor by hand]
Let $N=21$ and $a=2$. Compute the order $r$ of $2$ modulo $21$ by listing powers. Verify $r$ is even and that $2^{r/2}\not\equiv\pm1$, then compute $\gcd(2^{r/2}-1,21)$ and $\gcd(2^{r/2}+1,21)$ to factor $21$. Repeat with $a=5$ and report whether the lucky conditions hold.
\end{exercise}

\begin{exercise}[When Shor's luck fails]
Find an $a$ coprime to $N=15$ whose order modulo $15$ is even such that $a^{r/2} \equiv -1 \pmod{15}$ (Hint: try $a=14$), so that the order-finding fails to yield a non-trivial factor. Explain in one sentence why this condition blocks the construction, and what the algorithm does in that case.
\end{exercise}

\begin{exercise}[Grover arithmetic]
A cipher uses a $112$-bit key. Estimate the number of operations a classical brute force and an ideal Grover search each require, expressed as powers of two. If a designer wants $128$ bits of security against a quantum adversary using Grover, what is the smallest key length they should choose, and why is $256$ the usual round number?
\end{exercise}

\begin{exercise}[Why the error matters]
In the LWE setup, suppose all errors $e_i$ are exactly $0$. Given $n$ samples $(\mathbf a_i, b_i)$ with $\mathbf a_i\in\Zn{q}^n$, describe precisely how to recover the secret $\mathbf s$ in polynomial time, and state how large $n$ must be relative to the dimension. Then explain why adding even a small random error to each $b_i$ defeats your method.
\end{exercise}

\begin{exercise}[Decrypt a tiny LWE ciphertext]
Using the instance of Example~\ref{ch35:ex-lwe} ($q=17$, $\mathbf s=(4,1,11)$), an encryptor sends $\mathbf u=(2,2,2)$ and $v=14$ for some bit $\mu\in\{0,1\}$. Compute $v-\langle\mathbf s,\mathbf u\rangle \bmod 17$, fold the result into $[-8,8]$, and decide whether $\mu=0$ (result near $0$) or $\mu=1$ (result near $\floor{17/2}=8$).
\end{exercise}

\begin{exercise}[Break a two-time Lamport key]
A Lamport key signs $3$-bit digests, with secret pairs $(x_{i,0},x_{i,1})$ for $i=1,2,3$. The owner mistakenly signs digest $010$ and then digest $100$ with the same key. List exactly which secret values $x_{i,b}$ an attacker now holds, and determine which two additional digests (besides $010$ and $100$) the attacker can forge valid signatures for.
\end{exercise}

\begin{exercise}[Merkle authentication path]
A Merkle-tree signature scheme has $8$ leaves (one-time public keys $pk_0,\dots,pk_7$) and a tree of height $3$. The signer uses leaf $pk_5$. List, in order from leaf to root, exactly which sibling hashes form the authentication path, and write out the sequence of hash computations the verifier performs to recompute the root.
\end{exercise}

\begin{exercise}[Hybrid security argument]
A hybrid key exchange derives its session key as $K=\mathrm{KDF}(s_1\concat s_2)$, where $s_1$ comes from X25519 and $s_2$ from ML-KEM. Argue carefully that an adversary who can break X25519 (e.g.\ with a quantum computer) but not ML-KEM learns nothing about $K$, assuming the KDF is a good pseudorandom function. Then state the single condition under which the hybrid fails.
\end{exercise}

\begin{exercise}[Mosca's inequality]
Your organization stores medical records that must remain confidential for $25$ years. A full cryptographic migration is estimated to take $8$ years. Suppose expert consensus puts the arrival of a cryptographically relevant quantum computer at $20$ years out. Using $X+Y>Z$, determine whether you have a problem, and compute the latest year you could have started migration while still protecting all data through its confidentiality lifetime (take ``now'' as year $0$).
\end{exercise}

\begin{exercise}[Spot the structure]
For each problem, state whether Shor's algorithm solves it in polynomial time, and give a one-line reason based on whether it is an abelian hidden-period/subgroup problem: (a) factoring an RSA modulus; (b) elliptic-curve discrete log; (c) finding the shortest vector in a $400$-dimensional lattice; (d) inverting SHA-256; (e) computing a discrete log in $\Zns{p}$.
\end{exercise}

\begin{exercise}[Decryption failure as a leak]
Explain why a lattice KEM with a non-negligible decryption-failure probability is dangerous even if each individual failure seems harmless. Sketch how an attacker who can submit ciphertexts and observe whether decryption succeeded could, over many queries, learn information about the secret key. (This is why standardized parameters drive the failure rate below about $2^{-130}$.)
\end{exercise}
